\section{Results}

\subsection{Main Results}

Table~\ref{tab:main_results} presents the primary classification performance across three model stages: a pretrained baseline, a fine-tuned model, and a fine-tuned model augmented with synthetic dermoscopy images. Stages 2 and 3 share identical hyperparameters; the only difference is the addition of synthetic training images.

\begin{figure}[htbp]
  \centering
  \includegraphics[width=0.7\textwidth]{figures/fig2_subgroup_auroc.pdf}
  \caption{Subgroup AUROC across pipeline stages. Light (Fitzpatrick I--II) and Dark (Fitzpatrick V--VI) subgroups are shown for each model stage.}
  \label{fig:subgroup_auroc}
\end{figure}

Figure~\ref{fig:subgroup_auroc} visualizes the per-subgroup AUROC progression across these stages.

The pretrained baseline achieved an overall AUROC of 0.662 on the held-out test set. Stratified evaluation revealed a skin-tone disparity: the Dark subgroup AUROC was 0.588, while the Light subgroup AUROC was 0.719 (gap\,$=$\,$-0.131$). The baseline checkpoint was trained for up to 10 epochs with early stopping and reached validation AUROC 0.9650.

Fine-tuning on the DDI training split reduced overall AUROC to 0.634 (Table~\ref{tab:main_results}). Light subgroup AUROC was 0.756 and Dark subgroup AUROC was 0.547, widening the inter-group gap from $-0.131$ to $-0.209$. This result shows that standard fine-tuning did not remove the observed disparity.

The addition of synthetic augmentation yielded overall AUROC 0.640 and Dark subgroup AUROC 0.481. The Light--Dark gap widened to $-0.291$. The paired bootstrap comparison gave a Dark AUROC delta of $-0.066$ and $p=0.962$ (one-tailed directional, $H_1$: Dark improvement $>0$; $p$\,$=$\,proportion of bootstrap $\Delta \leq 0$). These findings indicate that simple photometric augmentation from 29 source images did not close the skin-tone performance gap.

\subsection{Subgroup Analysis}

Table~\ref{tab:subgroup_auroc} reports per-subgroup AUROC, sensitivity, and specificity across all three model stages.

\begin{figure}[htbp]
  \centering
  \includegraphics[width=0.7\textwidth]{figures/fig4_bootstrap_ci.pdf}
  \caption{Bootstrap 95\% confidence intervals for subgroup AUROC estimates (1000 iterations). Points indicate point estimates; bars indicate CI bounds.}
  \label{fig:bootstrap_ci}
\end{figure}

Figure~\ref{fig:bootstrap_ci} displays bootstrap confidence intervals for the subgroup AUROC estimates.

The Dark subgroup exhibited no consistent improvement across stages. AUROC progressed from 0.588 (baseline) to 0.547 (fine-tuned) to 0.481 (augmented), while Light AUROC progressed from 0.719 to 0.756 to 0.772.

Bootstrap confidence intervals (1000 iterations) confirmed the high uncertainty. The fine-tuned Dark AUROC bootstrap CI was $[0.341,\; 0.750]$, and the augmented model Dark AUROC CI was $[0.276,\; 0.688]$. Their overlap and the paired comparison p-value of 0.962 (one-tailed directional, $H_1$: Dark improvement $>0$) show that augmentation did not significantly improve Dark AUROC.

\subsection{Sensitivity Analysis}

Table~\ref{tab:sensitivity} reports sensitivity and specificity stratified by skin tone. Thresholds for all stages use Youden's $J$ on validation; AUROC is the primary comparison metric.

\begin{figure}[htbp]
  \centering
  \includegraphics[width=0.7\textwidth]{figures/fig3_sensitivity_gap.pdf}
  \caption{Sensitivity by skin tone across pipeline stages. Annotations show the Dark$-$Light sensitivity gap at each stage.}
  \label{fig:sensitivity_gap}
\end{figure}

Figure~\ref{fig:sensitivity_gap} illustrates the sensitivity gap trajectory across model stages.

Dark subgroup sensitivity was 0.400 (baseline), 0.600 (fine-tuned), and 0.500 (augmented). These threshold-dependent values do not establish a fairness improvement.

Light subgroup sensitivity was 0.700 (baseline), 0.600 (fine-tuned), and 0.600 (augmented).

Sensitivity values are computed over small melanoma-positive subgroups ($n=10$ per subgroup) and should be interpreted as directional rather than definitive.

\subsection{Ablation Study}

Table~\ref{tab:ablation} summarizes the ablation results. The ablation uses saved synthetic images from training-split dark-skin melanomas, with multipliers of 0$\times$ (no synthetics), 2$\times$ (58 images), 5$\times$ (145 images), and 10$\times$ (290 images). All ablation variants use identical hyperparameters to Stage 3.

\begin{figure}[htbp]
  \centering
  \includegraphics[width=0.7\textwidth]{figures/fig5_ablation_tradeoff.pdf}
  \caption{Ablation: validation AUROC, test AUROC, and dark-subgroup test AUROC across saved synthetic multipliers.}
  \label{fig:ablation_tradeoff}
\end{figure}

Figure~\ref{fig:ablation_tradeoff} presents the relationship between synthetic multiplier and performance across three metrics.

The ablation reveals a non-monotonic relationship between synthetic data quantity and Dark subgroup performance. Dark AUROC was 0.5438, 0.5750, 0.5344, and 0.4812 at 0$\times$, 2$\times$, 5$\times$, and 10$\times$, respectively. Overall test AUROC was 0.6848, 0.7063, 0.6742, and 0.6433, while validation AUROC peaked at 2$\times$. This exploratory pattern requires multi-seed confirmation.

Subgroup-level metrics for ablation variants show that the 2$\times$ variant had the highest Dark AUROC in this seed-42 run, while the 10$\times$ variant was lowest. The result should not be interpreted as a stable optimal multiplier without multi-seed repetition.
